% --- Archon physics formalization source begin ---
% archon:physics
% archon:covers PhyXMiniProblems/problem_phyx_mini_0398.lean
% archon:source-report reports/phyx_mini/problem_phyx_mini_0398.source.json

\chapter{Physics problem phyx\_mini\_0398}
\label{ch:PhyXMiniProblems_problem_phyx_mini_0398}

\paragraph{Problem source.}
\#\# Physical scenario

A pressure cooker has the lid screwed on tight. A small opening with A= 5 \$mm\textasciicircum{}2\$ is covered with a petcock that can be lifted to let steam escape.

\#\# Question

How much mass should the petcock have to allow boiling at 120◦C with an outside atmosphere at 101.3 kPa?

\#\# Answer choices

- A: A: 40g

- B: B: 50g

- C: C: 60g

- D: D: 70g

\#\# Figure caption (auxiliary; use the image as primary evidence)

The image illustrates a container with a liquid and vapor inside, possibly showing a phase change or steam generation process. The container is outlined with a solid line, and it has a dotted lid equipped with a valve at the top. 

Inside the container, the lower portion is labeled as "Liquid" and is represented with horizontal lines. The upper portion is labeled "Steam or vapor" and is filled with densely packed dots to suggest a gaseous state.

Above the lid, there is text labeled "Steam" with wavy lines indicating that steam is escaping from the valve. The overall diagram is likely meant to illustrate the concept of liquid turning into vapor within a closed system.

\#\# Dataset metadata

Laws of Thermodynamics; Physical Model Grounding Reasoning, Multi-Formula Reasoning

\paragraph{Recorded answer/context.}
B: B: 50g

\paragraph{Figure/image path.}
/root/proposal\_for\_physic/science-mango/phyx\_mini\_run/phyx\_data/test\_image/398.png

\paragraph{Formalization target.}
create a compiling Lean file with sorry bodies at `PhyXMiniProblems/problem\_phyx\_mini\_0398.lean`. The Lean declarations must preserve the physical quantities, dimensions or dimensional roles, figure labels, governing-law hypotheses, and final relation expressed by this problem.
Use Mathlib/Physlib names found through LeanExplore where available. If a physics API is missing, introduce faithful local abstractions rather than scalar placeholder aliases.

\begin{theorem}[Physics formalization target]
\label{thm:physics:phyx_mini_0398:target}
\lean{PhyXMiniProblems.ProblemPhyXMini0398.problem_phyx_mini_0398}
\uses{def:physics:phyx-mini-0398:phyxminiproblems-problemphyxmini0398-massingrams, def:physics:phyx-mini-0398:phyxminiproblems-problemphyxmini0398-pressurecookersetup, def:physics:phyx-mini-0398:phyxminiproblems-problemphyxmini0398-matchesprimaryfigure, def:physics:phyx-mini-0398:phyxminiproblems-problemphyxmini0398-matchesproblemdata, def:physics:phyx-mini-0398:phyxminiproblems-problemphyxmini0398-matchesstandardwaterdata, def:physics:phyx-mini-0398:phyxminiproblems-problemphyxmini0398-hasphysicalpetcockparameters, def:physics:phyx-mini-0398:phyxminiproblems-problemphyxmini0398-matchesidealpetcockmodel, def:physics:phyx-mini-0398:phyxminiproblems-problemphyxmini0398-satisfieswaterboilinglaw, def:physics:phyx-mini-0398:phyxminiproblems-problemphyxmini0398-atincipientventing, def:physics:phyx-mini-0398:phyxminiproblems-problemphyxmini0398-recordedanswerchoice, def:physics:phyx-mini-0398:phyxminiproblems-problemphyxmini0398-isuniqueclosestdisplayedmass}
The assigned autoformalize agent should translate this physics problem into Lean declarations in the covered file, with theorem and lemma proof bodies written as `by sorry`.
\end{theorem}
\begin{proof}
This is an autoformalization task, not a proof task. Produce statements that can later be proved by the physics prover without weakening the physical model.
\end{proof}

% --- Archon declaration topology begin ---
\section{Declaration topology}
\begin{definition}[Mass Quantity]
  \label{def:physics:phyx-mini-0398:phyxminiproblems-problemphyxmini0398-massquantity}
  \lean{PhyXMiniProblems.ProblemPhyXMini0398.MassQuantity}
  A nonnegative physical mass, independent of the unit used to read it
\end{definition}
\begin{proof}
  This is the declared model definition of Mass Quantity.
\end{proof}
\begin{definition}[Acceleration Quantity]
  \label{def:physics:phyx-mini-0398:phyxminiproblems-problemphyxmini0398-accelerationquantity}
  \lean{PhyXMiniProblems.ProblemPhyXMini0398.AccelerationQuantity}
  A nonnegative acceleration magnitude with physical dimension L T⁻²
\end{definition}
\begin{proof}
  This is the declared model definition of Acceleration Quantity.
\end{proof}
\begin{definition}[mass In Grams]
  \label{def:physics:phyx-mini-0398:phyxminiproblems-problemphyxmini0398-massingrams}
  \lean{PhyXMiniProblems.ProblemPhyXMini0398.massInGrams}
  \uses{def:physics:phyx-mini-0398:phyxminiproblems-problemphyxmini0398-massquantity}
  Read a physical mass in grams using Physlib's gram unit
\end{definition}
\begin{proof}
  This is the declared model definition of mass In Grams.
\end{proof}
\begin{definition}[area In Square Meters]
  \label{def:physics:phyx-mini-0398:phyxminiproblems-problemphyxmini0398-areainsquaremeters}
  \lean{PhyXMiniProblems.ProblemPhyXMini0398.areaInSquareMeters}
  Read the opening area in square metres
\end{definition}
\begin{proof}
  This is the declared model definition of area In Square Meters.
\end{proof}
\begin{definition}[area In Square Millimeters]
  \label{def:physics:phyx-mini-0398:phyxminiproblems-problemphyxmini0398-areainsquaremillimeters}
  \lean{PhyXMiniProblems.ProblemPhyXMini0398.areaInSquareMillimeters}
  Read the opening area in square millimetres
\end{definition}
\begin{proof}
  This is the declared model definition of area In Square Millimeters.
\end{proof}
\begin{definition}[acceleration In Meters Per Second Squared]
  \label{def:physics:phyx-mini-0398:phyxminiproblems-problemphyxmini0398-accelerationinmeterspersecondsquared}
  \lean{PhyXMiniProblems.ProblemPhyXMini0398.accelerationInMetersPerSecondSquared}
  \uses{def:physics:phyx-mini-0398:phyxminiproblems-problemphyxmini0398-accelerationquantity}
  Read an acceleration magnitude in metres per second squared
\end{definition}
\begin{proof}
  This is the declared model definition of acceleration In Meters Per Second Squared.
\end{proof}
\begin{definition}[pressure In Pascals]
  \label{def:physics:phyx-mini-0398:phyxminiproblems-problemphyxmini0398-pressureinpascals}
  \lean{PhyXMiniProblems.ProblemPhyXMini0398.pressureInPascals}
  Read an absolute pressure in pascals
\end{definition}
\begin{proof}
  This is the declared model definition of pressure In Pascals.
\end{proof}
\begin{definition}[pressure In Kilopascals]
  \label{def:physics:phyx-mini-0398:phyxminiproblems-problemphyxmini0398-pressureinkilopascals}
  \lean{PhyXMiniProblems.ProblemPhyXMini0398.pressureInKilopascals}
  \uses{def:physics:phyx-mini-0398:phyxminiproblems-problemphyxmini0398-pressureinpascals}
  Read an absolute pressure in kilopascals
\end{definition}
\begin{proof}
  This is the declared model definition of pressure In Kilopascals.
\end{proof}
\begin{definition}[temperature In Kelvin]
  \label{def:physics:phyx-mini-0398:phyxminiproblems-problemphyxmini0398-temperatureinkelvin}
  \lean{PhyXMiniProblems.ProblemPhyXMini0398.temperatureInKelvin}
  Read a Physlib absolute temperature in kelvin
\end{definition}
\begin{proof}
  This is the declared model definition of temperature In Kelvin.
\end{proof}
\begin{definition}[temperature In Degrees Celsius]
  \label{def:physics:phyx-mini-0398:phyxminiproblems-problemphyxmini0398-temperatureindegreescelsius}
  \lean{PhyXMiniProblems.ProblemPhyXMini0398.temperatureInDegreesCelsius}
  \uses{def:physics:phyx-mini-0398:phyxminiproblems-problemphyxmini0398-temperatureinkelvin}
  Convert the absolute-temperature readout to degrees Celsius
\end{definition}
\begin{proof}
  This is the declared model definition of temperature In Degrees Celsius.
\end{proof}
\begin{definition}[Cooker Fluid]
  \label{def:physics:phyx-mini-0398:phyxminiproblems-problemphyxmini0398-cookerfluid}
  \lean{PhyXMiniProblems.ProblemPhyXMini0398.CookerFluid}
  The working fluid whose vapor is labelled as steam in the figure
\end{definition}
\begin{proof}
  This is the declared model definition of Cooker Fluid.
\end{proof}
\begin{definition}[Lid Closure]
  \label{def:physics:phyx-mini-0398:phyxminiproblems-problemphyxmini0398-lidclosure}
  \lean{PhyXMiniProblems.ProblemPhyXMini0398.LidClosure}
  How the pressure-cooker lid is attached to the vessel
\end{definition}
\begin{proof}
  This is the declared model definition of Lid Closure.
\end{proof}
\begin{definition}[Opening Cover]
  \label{def:physics:phyx-mini-0398:phyxminiproblems-problemphyxmini0398-openingcover}
  \lean{PhyXMiniProblems.ProblemPhyXMini0398.OpeningCover}
  The device covering the small top opening
\end{definition}
\begin{proof}
  This is the declared model definition of Opening Cover.
\end{proof}
\begin{definition}[Figure Label]
  \label{def:physics:phyx-mini-0398:phyxminiproblems-problemphyxmini0398-figurelabel}
  \lean{PhyXMiniProblems.ProblemPhyXMini0398.FigureLabel}
  Text labels visibly printed in the supplied pressure-cooker figure
\end{definition}
\begin{proof}
  This is the declared model definition of Figure Label.
\end{proof}
\begin{definition}[Petcock Motion]
  \label{def:physics:phyx-mini-0398:phyxminiproblems-problemphyxmini0398-petcockmotion}
  \lean{PhyXMiniProblems.ProblemPhyXMini0398.PetcockMotion}
  Direction in which the petcock can move to uncover the opening
\end{definition}
\begin{proof}
  This is the declared model definition of Petcock Motion.
\end{proof}
\begin{definition}[Pressure Cooker Figure]
  \label{def:physics:phyx-mini-0398:phyxminiproblems-problemphyxmini0398-pressurecookerfigure}
  \lean{PhyXMiniProblems.ProblemPhyXMini0398.PressureCookerFigure}
  \uses{def:physics:phyx-mini-0398:phyxminiproblems-problemphyxmini0398-figurelabel}
  Qualitative information transcribed from image 398.png. The bitmap shows liquid below a steam-or-vapor region, a petcock at the top opening, and steam escaping above it. It contains no mass, pressure, area, or temperature value
\end{definition}
\begin{proof}
  This is the declared model definition of Pressure Cooker Figure.
\end{proof}
\begin{definition}[Pressure Cooker Setup]
  \label{def:physics:phyx-mini-0398:phyxminiproblems-problemphyxmini0398-pressurecookersetup}
  \lean{PhyXMiniProblems.ProblemPhyXMini0398.PressureCookerSetup}
  \uses{def:physics:phyx-mini-0398:phyxminiproblems-problemphyxmini0398-massquantity, def:physics:phyx-mini-0398:phyxminiproblems-problemphyxmini0398-accelerationquantity, def:physics:phyx-mini-0398:phyxminiproblems-problemphyxmini0398-cookerfluid, def:physics:phyx-mini-0398:phyxminiproblems-problemphyxmini0398-lidclosure, def:physics:phyx-mini-0398:phyxminiproblems-problemphyxmini0398-openingcover, def:physics:phyx-mini-0398:phyxminiproblems-problemphyxmini0398-petcockmotion, def:physics:phyx-mini-0398:phyxminiproblems-problemphyxmini0398-pressurecookerfigure}
  The pressure cooker and its independent physical observables. In particular, petcockMass is not defined from the recorded answer. The saturation-pressure function is retained as thermodynamic data rather than replacing pressure by a bare scalar
\end{definition}
\begin{proof}
  This is the declared model definition of Pressure Cooker Setup.
\end{proof}
\begin{definition}[Matches Primary Figure]
  \label{def:physics:phyx-mini-0398:phyxminiproblems-problemphyxmini0398-matchesprimaryfigure}
  \lean{PhyXMiniProblems.ProblemPhyXMini0398.MatchesPrimaryFigure}
  \uses{def:physics:phyx-mini-0398:phyxminiproblems-problemphyxmini0398-pressurecookersetup}
  Qualitative labels and geometry read directly from the primary image
\end{definition}
\begin{proof}
  This is the declared model definition of Matches Primary Figure.
\end{proof}
\begin{definition}[Matches Problem Data]
  \label{def:physics:phyx-mini-0398:phyxminiproblems-problemphyxmini0398-matchesproblemdata}
  \lean{PhyXMiniProblems.ProblemPhyXMini0398.MatchesProblemData}
  \uses{def:physics:phyx-mini-0398:phyxminiproblems-problemphyxmini0398-areainsquaremillimeters, def:physics:phyx-mini-0398:phyxminiproblems-problemphyxmini0398-pressureinkilopascals, def:physics:phyx-mini-0398:phyxminiproblems-problemphyxmini0398-temperatureindegreescelsius, def:physics:phyx-mini-0398:phyxminiproblems-problemphyxmini0398-pressurecookersetup}
  The apparatus description and numerical data stated in the problem. No mass value or answer-choice label occurs here
\end{definition}
\begin{proof}
  This is the declared model definition of Matches Problem Data.
\end{proof}
\begin{definition}[Matches Standard Water Data]
  \label{def:physics:phyx-mini-0398:phyxminiproblems-problemphyxmini0398-matchesstandardwaterdata}
  \lean{PhyXMiniProblems.ProblemPhyXMini0398.MatchesStandardWaterData}
  \uses{def:physics:phyx-mini-0398:phyxminiproblems-problemphyxmini0398-accelerationinmeterspersecondsquared, def:physics:phyx-mini-0398:phyxminiproblems-problemphyxmini0398-pressureinkilopascals, def:physics:phyx-mini-0398:phyxminiproblems-problemphyxmini0398-pressurecookersetup}
  Standard terrestrial gravity and a rounded water steam-table readout at 120 °C. The saturation-pressure value is empirical input to the model; it does not state the requested petcock mass
\end{definition}
\begin{proof}
  This is the declared model definition of Matches Standard Water Data.
\end{proof}
\begin{definition}[Has Physical Petcock Parameters]
  \label{def:physics:phyx-mini-0398:phyxminiproblems-problemphyxmini0398-hasphysicalpetcockparameters}
  \lean{PhyXMiniProblems.ProblemPhyXMini0398.HasPhysicalPetcockParameters}
  \uses{def:physics:phyx-mini-0398:phyxminiproblems-problemphyxmini0398-massingrams, def:physics:phyx-mini-0398:phyxminiproblems-problemphyxmini0398-areainsquaremeters, def:physics:phyx-mini-0398:phyxminiproblems-problemphyxmini0398-accelerationinmeterspersecondsquared, def:physics:phyx-mini-0398:phyxminiproblems-problemphyxmini0398-pressureinpascals, def:physics:phyx-mini-0398:phyxminiproblems-problemphyxmini0398-pressurecookersetup}
  Positivity and pressure ordering for the intended operating state
\end{definition}
\begin{proof}
  This is the declared model definition of Has Physical Petcock Parameters.
\end{proof}
\begin{definition}[Matches Ideal Petcock Model]
  \label{def:physics:phyx-mini-0398:phyxminiproblems-problemphyxmini0398-matchesidealpetcockmodel}
  \lean{PhyXMiniProblems.ProblemPhyXMini0398.MatchesIdealPetcockModel}
  \uses{def:physics:phyx-mini-0398:phyxminiproblems-problemphyxmini0398-pressurecookersetup}
  Idealizations used by the one-dimensional petcock force balance
\end{definition}
\begin{proof}
  This is the declared model definition of Matches Ideal Petcock Model.
\end{proof}
\begin{definition}[Satisfies Water Boiling Law]
  \label{def:physics:phyx-mini-0398:phyxminiproblems-problemphyxmini0398-satisfieswaterboilinglaw}
  \lean{PhyXMiniProblems.ProblemPhyXMini0398.SatisfiesWaterBoilingLaw}
  \uses{def:physics:phyx-mini-0398:phyxminiproblems-problemphyxmini0398-pressurecookersetup}
  At equilibrium boiling, the cooker steam pressure equals the water saturation pressure at the water temperature. This generic law contains no numerical mass or answer choice
\end{definition}
\begin{proof}
  This is the declared model definition of Satisfies Water Boiling Law.
\end{proof}
\begin{definition}[upward Steam Pressure Force In Newtons]
  \label{def:physics:phyx-mini-0398:phyxminiproblems-problemphyxmini0398-upwardsteampressureforceinnewtons}
  \lean{PhyXMiniProblems.ProblemPhyXMini0398.upwardSteamPressureForceInNewtons}
  \uses{def:physics:phyx-mini-0398:phyxminiproblems-problemphyxmini0398-areainsquaremeters, def:physics:phyx-mini-0398:phyxminiproblems-problemphyxmini0398-pressureinpascals, def:physics:phyx-mini-0398:phyxminiproblems-problemphyxmini0398-pressurecookersetup}
  Upward steam-pressure force on the underside of the petcock, in newtons
\end{definition}
\begin{proof}
  This is the declared model definition of upward Steam Pressure Force In Newtons.
\end{proof}
\begin{definition}[downward Atmospheric Force In Newtons]
  \label{def:physics:phyx-mini-0398:phyxminiproblems-problemphyxmini0398-downwardatmosphericforceinnewtons}
  \lean{PhyXMiniProblems.ProblemPhyXMini0398.downwardAtmosphericForceInNewtons}
  \uses{def:physics:phyx-mini-0398:phyxminiproblems-problemphyxmini0398-areainsquaremeters, def:physics:phyx-mini-0398:phyxminiproblems-problemphyxmini0398-pressureinpascals, def:physics:phyx-mini-0398:phyxminiproblems-problemphyxmini0398-pressurecookersetup}
  Downward outside-atmosphere force on the top face, in newtons
\end{definition}
\begin{proof}
  This is the declared model definition of downward Atmospheric Force In Newtons.
\end{proof}
\begin{definition}[petcock Weight In Newtons]
  \label{def:physics:phyx-mini-0398:phyxminiproblems-problemphyxmini0398-petcockweightinnewtons}
  \lean{PhyXMiniProblems.ProblemPhyXMini0398.petcockWeightInNewtons}
  \uses{def:physics:phyx-mini-0398:phyxminiproblems-problemphyxmini0398-massingrams, def:physics:phyx-mini-0398:phyxminiproblems-problemphyxmini0398-accelerationinmeterspersecondsquared, def:physics:phyx-mini-0398:phyxminiproblems-problemphyxmini0398-pressurecookersetup}
  Petcock weight in newtons, using grams for the displayed mass readout
\end{definition}
\begin{proof}
  This is the declared model definition of petcock Weight In Newtons.
\end{proof}
\begin{definition}[At Incipient Venting]
  \label{def:physics:phyx-mini-0398:phyxminiproblems-problemphyxmini0398-atincipientventing}
  \lean{PhyXMiniProblems.ProblemPhyXMini0398.AtIncipientVenting}
  \uses{def:physics:phyx-mini-0398:phyxminiproblems-problemphyxmini0398-pressurecookersetup, def:physics:phyx-mini-0398:phyxminiproblems-problemphyxmini0398-upwardsteampressureforceinnewtons, def:physics:phyx-mini-0398:phyxminiproblems-problemphyxmini0398-downwardatmosphericforceinnewtons, def:physics:phyx-mini-0398:phyxminiproblems-problemphyxmini0398-petcockweightinnewtons}
  Vertical force balance at incipient lift: the steam force just balances the outside atmospheric force together with the petcock weight. It is a generic governing relation and does not fix the unknown mass
\end{definition}
\begin{proof}
  This is the declared model definition of At Incipient Venting.
\end{proof}
\begin{definition}[Answer Choice]
  \label{def:physics:phyx-mini-0398:phyxminiproblems-problemphyxmini0398-answerchoice}
  \lean{PhyXMiniProblems.ProblemPhyXMini0398.AnswerChoice}
  Labels of the four answers printed with the problem
\end{definition}
\begin{proof}
  This is the declared model definition of Answer Choice.
\end{proof}
\begin{definition}[displayed Mass In Grams]
  \label{def:physics:phyx-mini-0398:phyxminiproblems-problemphyxmini0398-displayedmassingrams}
  \lean{PhyXMiniProblems.ProblemPhyXMini0398.displayedMassInGrams}
  \uses{def:physics:phyx-mini-0398:phyxminiproblems-problemphyxmini0398-answerchoice}
  Petcock mass in grams printed beside each answer label
\end{definition}
\begin{proof}
  This is the declared model definition of displayed Mass In Grams.
\end{proof}
\begin{definition}[recorded Answer Choice]
  \label{def:physics:phyx-mini-0398:phyxminiproblems-problemphyxmini0398-recordedanswerchoice}
  \lean{PhyXMiniProblems.ProblemPhyXMini0398.recordedAnswerChoice}
  \uses{def:physics:phyx-mini-0398:phyxminiproblems-problemphyxmini0398-answerchoice}
  The answer label recorded in the supplied dataset
\end{definition}
\begin{proof}
  This is the declared model definition of recorded Answer Choice.
\end{proof}
\begin{definition}[Is Unique Closest Displayed Mass]
  \label{def:physics:phyx-mini-0398:phyxminiproblems-problemphyxmini0398-isuniqueclosestdisplayedmass}
  \lean{PhyXMiniProblems.ProblemPhyXMini0398.IsUniqueClosestDisplayedMass}
  \uses{def:physics:phyx-mini-0398:phyxminiproblems-problemphyxmini0398-answerchoice, def:physics:phyx-mini-0398:phyxminiproblems-problemphyxmini0398-displayedmassingrams}
  A choice is the unique closest displayed mass to an unrounded result
\end{definition}
\begin{proof}
  This is the declared model definition of Is Unique Closest Displayed Mass.
\end{proof}
% --- Archon declaration topology end ---
% --- Archon physics formalization source end ---
